\documentclass[12pt,letterpaper]{article}



%%
%% Required packages:
%%
\usepackage{etex}
\usepackage{amsmath}
\usepackage{amssymb}
\usepackage{amstext}
\usepackage{amsthm}
\usepackage{fancyhdr,fancybox}
\usepackage{mathrsfs} % need for \mathscr command
\usepackage{multicol}
\usepackage[normalem]{ulem}
%\usepackage{pictex,pictexplus}
% \usepackage{pictexwd}
\usepackage{rotating}
\usepackage{url}
\usepackage{xfrac}
\usepackage{booktabs}
\usepackage{color,soul}
\usepackage{comment}

\usepackage{hyperref}
\hypersetup{
    colorlinks=true,     % false: boxed links; true: colored links
    linkcolor=blue,      % color of internal links
    citecolor=blue,      % color of links to bibliography
    filecolor=blue,      % color of file links
    urlcolor=blue        % color of external links
}


\usepackage{tikz}
\usetikzlibrary{backgrounds,calc}

%%
%% Common formatting commands
%%
\setlength{\parindent}{0in}
\setlength{\textwidth}{7in}
\setlength{\evensidemargin}{-0.25in}
\setlength{\oddsidemargin}{-0.25in}
\setlength{\parskip}{.5\baselineskip}
\setlength{\topmargin}{-0.5in}
\setlength{\textheight}{9in}

%               Problem and Part
\newcounter{problemnumber}
\newcounter{partnumber}
\newcounter{subpartnumber}
\newcommand{\Problem}{%
  \refstepcounter{problemnumber}%
  \setcounter{partnumber}{0}%
  \setcounter{subpartnumber}{0}%
  \item[\makebox{\hfil\textbf{\theproblemnumber.}\hfil}]%
  }
\newcommand{\InProblem}[2][1.6]{%
  \refstepcounter{problemnumber}%
  \setcounter{partnumber}{0}%
  \setcounter{subpartnumber}{0}%
  \textbf{\theproblemnumber.}\ \ \parbox[t]{#1in}{#2}%
}
\newcommand{\InSmallProblem}[1]{\InProblem[1.05]{#1}}
\newcommand{\NoProblem}[1]{%
  \mbox{\hfil\textbf{#1}\hfil}%
  }
\newcommand{\UnProblem}{%
  \refstepcounter{problemnumber}%
  \setcounter{partnumber}{0}%
  \setcounter{subpartnumber}{0}%
  \mbox{\hfil\textbf{\theproblemnumber.}\hfil}%
  }
\newcommand{\InFlexProblem}[1]{%
  \refstepcounter{problemnumber}%
  \setcounter{partnumber}{0}%
  \setcounter{subpartnumber}{0}%
  \textbf{\theproblemnumber.}\ \ #1%
}
%%  Part:
\newcommand{\Part}{%
  \renewcommand{\thepartnumber}{(\alph{partnumber})}%
  \refstepcounter{partnumber}
  \setcounter{subpartnumber}{0}%
  \item[(\alph{partnumber})]%
}
\newcommand{\InPart}[2][1.75]{%
  \renewcommand{\thepartnumber}{(\alph{partnumber})}%
  \refstepcounter{partnumber}%
  \setcounter{subpartnumber}{0}%
  (\alph{partnumber})\ \ \parbox[t]{#1in}{#2}%
}
\newcommand{\UnPart}{%
  \refstepcounter{partnumber}%
  \renewcommand{\thepartnumber}{(\alph{partnumber})}%
  \setcounter{subpartnumber}{0}%
  (\alph{partnumber})%
}
\newcommand{\InLargePart}[1]{\InPart[3.05]{#1}}
\newcommand{\InFullPart}[1]{\InPart[6.2]{#1}}
\newcommand{\InSmallPart}[1]{\InPart[1.05]{#1}}
%% SubPart:
\newcommand{\SubPart}{%
  \refstepcounter{subpartnumber}%
  \item[(\roman{subpartnumber})]%
  }
\newcommand{\InSubPart}[2][2.25]{%
  \stepcounter{subpartnumber}%
  (\roman{subpartnumber})\ \ \parbox[t]{#1in}{#2}%
}
\newcommand{\InSmallSubPart}[1]{\InSubPart[1.5]{#1}}
\newcommand{\InTinySubPart}[1]{%
  \stepcounter{subpartnumber}%
  (\roman{subpartnumber})\ \ \makebox{#1}%
}
\newcommand{\InMySubPart}[2][2.25]{%
  \stepcounter{subpartnumber}%
  \parbox[t]{#1in}{\makebox[0.3in][l]{(\roman{subpartnumber})}\ \ {#2}}%
}
\newcommand{\TrueFalse}{\textbf{T}\ \ \textbf{F}\ \ }
\newcommand{\TFPart}[1]{% 
  \stepcounter{partnumber}%
  \setcounter{subpartnumber}{0}%
  \item[(\alph{partnumber})] \TrueFalse\parbox[t]{5.5in}{#1}}

\newcommand{\Points}[1]{(#1 points) \ }
\newcommand{\Point}[1]{(#1 point) \ }


%% For Answers:
\newcounter{answernumber}
\newcommand{\Answer}{%
  \refstepcounter{answernumber}%
  \setcounter{partnumber}{0}%
  \setcounter{subpartnumber}{0}%
  \item[\makebox{\hfil\textbf{\theanswernumber.}\hfil}]%
  }


% Setting up the headers for the first page
%\chead{} % will default to what is typed in the file.
\fancyhead[L]{\textsc{Math 22a}}
\fancyhead[R]{\textsc{Fall 2024}}
\fancyfoot[R]{
	\vspace*{\fill}\footnotesize\textcopyright{} \emph{Harvard Math 22a}	
}
\renewcommand{\headrulewidth}{0pt}

%\fancyhead[C]{}
%	\chead{}	
%	\lhead{}
%	\rhead{}
%	\cfoot{}


% Putting copyright on the footer of each page
%\fancypagestyle{secondpage}{
%	\fancyhead[C]{TESTing again} % change this to be blank
%		%\chead{}	
%	\fancyhead[L]{bears} % change this to be blank
%	\fancyhead[R]{dogs} % change this to be blank
%	\fancyfoot[R]{ %keeps the footer the same
%		\vspace*{\fill}\footnotesize\textcopyright{} \emph{Harvard Math 22a}	
%	}
%}
%\renewcommand{\headrulewidth}{0pt}
%\pagestyle{fancy}
\pagestyle{empty}


%\lhead{}
%\rhead{}
%\rfoot{\vspace*{\fill}\footnotesize\textcopyright{} \emph{Harvard Math 22a}}
%\renewcommand{\headrulewidth}{0pt}
%\pagestyle{fancy}




%%
%% Common shortcut commands:
%%
\newcommand{\ds}{\displaystyle}
\newcommand{\sgn}{\operatorname{sgn}}
\renewcommand{\Pr}{\operatorname{P}}
\newcommand{\CPr}[3][{}]{\Pr_{\text{#1}}\left(#2\,|\,#3\right)}

\newcommand{\C}{\mathbb{C}}
\newcommand{\R}{\mathbb{R}}
\newcommand{\Z}{\mathbb{Z}}
\newcommand{\N}{\mathbb{N}}
\newcommand{\Q}{\mathbb{Q}}
\newcommand{\D}{\mathbb{D}}

\newcommand{\rref}{\operatorname{rref}}
\newcommand{\im}{\operatorname{im}}
\newcommand{\rank}{\operatorname{rank}}
\newcommand{\diag}{\operatorname{diag}}
\newcommand{\Span}{\operatorname{span}}
%\renewcommand{\vec}[1]{\mathbf{#1}}
\newcommand{\charpoly}{\mathscr{P}}
\newcommand{\nicefrac}[2]{\sfrac{#1}{#2}} % using xfrac package
\newcommand{\topology}{\mathcal{T}}
\newcommand{\Inn}{\operatorname{Inn}}

\newcommand{\blankbox}[2]{\fbox{\rule{#1}{0in}\rule{0in}{#2}}}

\newenvironment{myindentpar}[1]%
 {\begin{list}{}%
         {\setlength{\leftmargin}{#1}
          \setlength{\rightmargin}{\leftmargin}}%
         \item[]%
 }
 {\end{list}}
\newtheorem{theorem}{Theorem}
\newtheorem*{NStheorem}{Theorem}
\newtheorem{cor}[theorem]{Corollary}
\newtheorem*{claim}{Claim}
\newtheorem{defn}{Definition}

\newenvironment{amatrix}[1]{%
  \left(\begin{array}{@{}*{#1}{c}|c@{}}
}{%
  \end{array}\right)
}

%--------grstep
% For denoting a Gauss' reduction step.
% Use as: \grstep{\rho_1+\rho_3} or \grstep[2\rho_5 \\ 3\rho_6]{\rho_1+\rho_3}
\newcommand{\grstep}[2][\relax]{%
   \ensuremath{\mathrel{
       {\mathop{\longrightarrow}\limits^{#2\mathstrut}_{
                                     \begin{subarray}{l} #1 \end{subarray}}}}}}
\newcommand{\swap}{\leftrightarrow}


\usetikzlibrary{arrows}
\usepackage{wrapfig}
\usepackage{amsfonts,amsmath}

\makeatletter
\renewcommand*\env@matrix[1][*\c@MaxMatrixCols c]{%
	\hskip -\arraycolsep
	\let\@ifnextchar\new@ifnextchar
	\array{#1}}
\makeatother

\def\env@matrix{\hskip -\arraycolsep
	\let\@ifnextchar\new@ifnextchar
	\array{*\c@MaxMatrixCols c}}

\newcommand{\norm}[1]{\left\lVert#1\right\rVert}

\chead{\Large\textbf{Orthogonal Sets, $\vec\S$6.2, 6.3}}% 

\begin{document}
\thispagestyle{fancy}

Recall from last class:

\begin{itemize}
\item Let $u, v \in \mathbb{R}^n$, then the {\bf dot product} of $u$ and $v$ is given by $$u \cdot v= u^{T}v=u_1 v_1+\dots+ u_n v_n.$$
\item {\bf Theorem 6.1:} Let $u, v, w \in \mathbb{R}^n$ and $c\in \mathbb{R}$.
\begin{enumerate}
	\item $$ u \cdot v = v \cdot u.$$
	\item $$(u+v)\cdot w= u \cdot w+ v \cdot w.$$
	\item $$(cu)\cdot v = c(u\cdot v)= u \cdot (cv).$$
	\item $u \cdot u \geq 0$ and $u \cdot u =0$ if and only if $u=0$.
\end{enumerate}
\item  The {\bf length} or {\bf norm} of the vector $v$ is a non-negative scalar, denoted by $\norm{v}$, and defined by $$\norm{v}= \sqrt{v\cdot v}.$$
\item We define the {\bf distance} between $u$ and $v$ by $$\text{dist}(u,v)=\norm{u-v}.$$
\item The vectors $u$ and $v$ are {\bf orthogonal} if $u\cdot v=0$.
\item {\bf Theorem 6.2:} The vectors $u$ and $v$ are orthongal if and only if $$\norm{u}^2+\norm{v}^2=\norm{u+v}^2.$$
\item The set $$W^{\perp} =\{ z \in \mathbb{R}^n: z \cdot w=0 \text{ for all } w \in W\}$$ is called the {\bf orthogonal complement} of $W$.
\item Note: $x \in W^{\perp}$ if and only if $x$ is orthogonal to each vector in a set that spans $W$. Also $W^{\perp}$ is a subspace of $\mathbb{R}^n$.
\item We can use the dot product to find the angles between vectors; namely $$u\cdot v = \norm{u} \norm{v} \cos \theta.$$
\item 	A set of vectors $\{ u_1, \dots, u_p\}$ in $\mathbb{R}^n$ is {\bf orthogonal} if $u_k \cdot u_j=0$ for all $j\neq k$.
\item {\bf Theorem 6.3} Let $A$ be an $m\times n$ matrix, then $$(\text{Row}(A))^{\perp} = \text{Null}(A) \text{ and } (\text{Col}(A))^{\perp}= \text{Null}(A^{T}).$$
\item
{\bf Theorem 6.4:} If $S=\{ u_1, \dots, u_p\}$ is an orthogonal set of non-zero vectors in $\mathbb{R}^n$, then $S$ is linearly independent and hence a basis for $\text{Span} \; S$.
\end{itemize}

\newpage

\begin{defn} An {\bf orthogonal basis} for a subspace $W$ of $\mathbb{R}^n$ is a basis for $W$ that is also an orthogonal set.
\end{defn}

{\bf Theorem 6.5:} Let $\{u_1, \dots,u_p\}$ be an orthogonal basis for a subspace $W$ of $\mathbb{R}^n$. For each $y\in W$, then $$y=c_1 u_1+\dots+c_p u_p$$ where $$c_j=\frac{y\cdot u_j}{u_j \cdot u_j}.$$

\begin{enumerate}

\Problem Prove the Theorem.





\vfill


\Problem Let $u_1=\begin{bmatrix} 2\\ -3\\ \end{bmatrix}$ and $u_2= \begin{bmatrix} 6\\4\\ \end{bmatrix}$ and $x=\begin{bmatrix} 9\\ -7\\ \end{bmatrix}$. Find $x$ in the $\{u_1, u_2\}$-coordinates.

\vfill


\newpage

\begin{defn} Let $u$ be a non-zero vector in $\mathbb{R}^n$ and let $v$ be a vector in $\mathbb{R}^n$, then the {\bf orthogonal projection} of $v$ onto $u$, denoted by $\text{proj}_{u}{v}$, is given by $$\text{proj}_{u}{v}=\frac{v\cdot u}{u \cdot u} u$$ Let $L=\text{Span}\{ u\}$, then we also write $\text{proj}_{L}{v}$ for the orthogonal projection of $v$ onto the line defined by the span of $u$; namely $$\text{proj}_{L}{v}=\text{proj}_{u}{v}.$$
\end{defn}

\Problem Let $u=\begin{bmatrix} 2\\ 2\\ \end{bmatrix}$ and $y= \begin{bmatrix} 1\\4\\ \end{bmatrix}$. Find the orthogonal projection of $y$ onto $u$.

\vfill

{\bf Note:} Theorem 6.5 says that $$y=\text{proj}_{u_1}{y}+\text{proj}_{u_2}{y}+\dots+ \text{proj}_{u_p}{y}.$$

\begin{defn} A set $\{u_1, \dots, u_p\}$ is {\bf orthonormal} if it is orthogonal and $\norm{u_k}=1$ for all $k$.
\end{defn}

{\bf Theorem 6.6:} An $m\times n$ matrix $U$ has orthonormal columns if and only if $U^{T}U=I_n$.

\Problem Prove the Theorem.

\vfill

\newpage

{\bf Theorem 6.7:} Let $U$ be an $m\times n$ matrix with orthonormal columns. Let $x,y \in \mathbb{R}^n$. Then
\begin{enumerate}
\item $\norm{Ux}=\norm{x}$.
\item $(Ux)\cdot (Uy) = x \cdot y$.
\item $(Ux) \cdot (Uy)=0$ if and only if $x \cdot y=0$.
\end{enumerate}
\Problem Prove the theorem.

\vfill

\begin{defn} An {\bf orthogonal matrix} $U$ is a square matrix so that $U^{-1}=U^{T}$. 
\end{defn}

{\bf Note:} An $n \times n$ matrix $U$ is orthogonal if and only if the columns of $U$ are orthonormal.

\Problem Suppose $U$ is an orthogonal matrix.
\begin{enumerate}
\item  What can you say about the eigenvalues of $U$?

\vfill

\item What is the determinant of $U$?

\vfill

\end{enumerate}

\newpage



\end{enumerate}

\begin{comment}
\newpage
\chead{\Large\textbf{Orthogonal Sets -- Answers / Solutions}}%
\lhead{}
\rhead{}
\thispagestyle{fancy}

\begin{enumerate}
\Answer
\begin{proof} Suppose $y= c_1 u_1+\dots c_pu_p$. Now take the dot product of both sides with $u_k$ to get $$y\cdot u_k = (c_1 u_1+\dots +c_p u_p)\cdot u_k.$$ Using the orthogonality of $u_1, \dots, u_p$, we can simplify the right hand side to get $$y \cdot u_k= c_k u_k \cdot u_k.$$ Therefore $$c_k=\frac{y \cdot u_k}{u_k\cdot u_k}.$$
\end{proof}
\Answer In order to find $x$ in the $u_1, u_2$ coordinates, we need to find scalars $c_1, c_2$ such that $$x=c_1 u_1+ c_2 u_2.$$ Notice that $u_1\cdot u_2=0$, and hence we can use the previous theorem to conclude $$c_1=\frac{x\cdot u_1}{u_1\cdot u_1}= \frac{39}{13}=3,$$ and $$c_2=\frac{x\cdot u_2}{u_2\cdot u_2}= \frac{26}{52}=\frac{1}{2}.$$ Thus $[x]_{\{u_1,u_1\}}=\begin{bmatrix} 3\\ \frac{1}{2}\\ \end{bmatrix}$.
 
\Answer The orthogonal projection of $y$ onto $u$ is given by $$\text{proj}_{u}{y}=\frac{y \cdot u}{u \cdot u} u= \frac{10}{8}u=\begin{bmatrix}\frac{5}{2}\\[0.5em] \frac{5}{2}\\ \end{bmatrix}.$$

\Answer

\begin{proof}
	Let $U=[u_1, \dots, u_n]$, then $$U^{T}U=\begin{bmatrix} u_1^{T}\\ \vdots\\ u_n^{T} \end{bmatrix} \begin{bmatrix} u_1 & \dots & u_n\\ \end{bmatrix}=\begin{bmatrix}
	u_1^{T}u_1 & u_1^{T}u_2 & \dots & u_1^{T}u_n\\
	u_2^{T}u_1 & u_2^{T}u_2 & \dots & u_2^{T}u_n\\
	\vdots     &            & \ddots& \\
	u_n^{T}u_1 & u_n^{T}u_2 & \dots & u_n^{T}u_n\\ \end{bmatrix}=\begin{bmatrix}
	u_1 \cdot u_1 & u_1 \cdot u_2 & \dots & u_1 \cdot u_n\\
	u_2 \cdot u_1 & u_2 \cdot u_2 & \dots & u_2 \cdot u_n\\
	\vdots     &            & \ddots& \\
	u_n \cdot u_1 & u_n \cdot u_2 & \dots & u_n \cdot u_n\\ \end{bmatrix}.$$ Thus $U^{T}U=I_n$ if and only if $u_k \cdot u_j =0$ for all $k \neq j$ and $u_k \cdot u_k =1$. Thus $U^{T}U=I_n$ if and only if the columns of $U$ are orthonormal. 
\end{proof}
\Answer 
\begin{proof} Let $U$ be an $m\times n$ matrix with orthonormal columns. We prove (b) first.
	
{\bf (b).} By the previous theorem we know that $U^{T}U=I_n$. Now using the definition of the dot product and properties of the transpose of a matrix, we get  $$(Ux)\cdot (Uy)= (Ux)^{T}Uy = x^{T}U^{T} Uy =x^{T}y= x \cdot y.$$

{\bf (a).} By definition of the norm, $\norm{Ux}^2 = (Ux)\cdot (Ux)$. From part (a), we get $\norm{Ux}^2 = x \cdot x =\norm{x}^2$. Taking square roots, we get the result.

{\bf (c).} From part (b), $x \cdot y=0$ if and only if $(Ux)\cdot (Uy)=0$.
\end{proof}
\Answer
\begin{enumerate}
	\item Suppose $\lambda$ is an eigenvalue of an orthogonal matrix $U$. Since orthogonal matrices have orthonormal columns, we know that $U^{T}U=I_n$. Thus $$\norm{x}=\norm{U x}=\norm{\lambda x} =|\lambda| \norm{x}.$$ Hence we get $|\lambda|=1$.
	\item Using the multiplicativity of the determinant we get $$1=\det{I_n} = \det{U^T U}=\det{U^{T}}\det{U}= (\det{U})^2.$$ Thus $\det{U}=1$ or $\det{U}=-1$.
\end{enumerate}
\end{enumerate}  

\end{comment}

\end{document}


%%% Local ables: 
%%% mode: latex
%%% TeX-master: t
%%% End: 
